\section*{अध्याय \ref{ch:b1:acetals-protection} --- कार्बोनिल: नाभिकरागी योग, ऐसीटैल और रक्षण}

\begin{solution}{exo:b1:acetals-protection:1}
\ce{CH3CH(OH)OCH3} एक \omterm{def:b1:acetals-protection:hemiacetal}{हेमीऐसीटैल} है। \ce{CH3CH(OCH3)2} एक \omterm{def:b1:acetals-protection:hemiacetal}{ऐसीटैल} है।
\ce{CH3CH(OH)2} एक हाइड्रेट है। \ce{CH3OCH3} एक ईथर है। \ce{(CH3)2C(OCH2CH3)2}
एक \omterm{def:b1:acetals-protection:hemiacetal}{ऐसीटैल} है (एक कीटोन का)।
\end{solution}

\begin{solution}{exo:b1:acetals-protection:2}
\ce{CH3CH2CHO + 2CH3OH <=> CH3CH2CH(OCH3)2 + H2O}; और एथेन-1,2-डाइऑल के साथ
प्रोपेनोन, 2,2-डाइमेथिल-1,3-डाइऑक्सोलेन और जल देता है,
\ce{(CH3)2CO + HOCH2CH2OH <=> C5H10O2 + H2O}।
\end{solution}

\begin{solution}{exo:b1:acetals-protection:3}
सोडियम हाइड्रॉक्साइड, मेथिलमैग्नीशियम ब्रोमाइड और सोडियम बोरोहाइड्राइड उसे
अपरिवर्तित छोड़ देते हैं; पर तनु जलीय सल्फ़्यूरिक अम्ल उसका जल-अपघटन करता है।
\end{solution}

\begin{solution}{exo:b1:acetals-protection:4}
कार्बन 4 का OH ऐल्डिहाइड कार्बन 1 से जुड़ता है: चार कार्बनों और एक ऑक्सीजन का
पाँच-सदस्यीय वलय, जिसमें वलय-ऑक्सीजन के बगल वाले कार्बन पर एक OH
है।
\end{solution}

\begin{solution}{exo:b1:acetals-protection:5}
C=O का प्रोटॉनन; मेथेनॉल कार्बन से जुड़ता है; \ce{H+} की हानि: \omterm{def:b1:acetals-protection:hemiacetal}{हेमीऐसीटैल};
उसके OH का प्रोटॉनन; जल की हानि से \ce{CH3CH=O+CH3}; दूसरा मेथेनॉल
जुड़ता है; \ce{H+} की हानि: \ce{CH3CH(OCH3)2}। लिए गए दोनों प्रोटॉन
लौटा दिए जाते हैं: अम्ल एक \omterm{def:b1:elementary-steps:catalyst}{उत्प्रेरक} है।
\end{solution}

\begin{solution}{exo:b1:acetals-protection:6}
एक OCH3 का प्रोटॉनन; मेथेनॉल की हानि से \ce{(CH3)2C=O+CH3}; जल जुड़ता है;
\ce{H+} की हानि: \omterm{def:b1:acetals-protection:hemiacetal}{हेमीऐसीटैल}; प्रोटॉनन, मेथेनॉल की हानि, \ce{H+} की हानि:
प्रोपेनोन। आधिक्य जल जल-अपघटन के लिए $Q$ को $K$ से नीचे रखता है, और वह
पूर्णता तक चलता है।
\end{solution}

\begin{solution}{exo:b1:acetals-protection:7}
प्रति \omterm{def:b1:acetals-protection:hemiacetal}{ऐसीटैल} एक जल: $0.250 \times 18.0 = \qty{4.5}{g}$, लगभग
\qty{4.5}{mL}। एक घंटे बाद, $3.2/18.0 = \qty{0.178}{mol}$: \qty{71}{\percent}
रूपांतरण।
\end{solution}

\begin{solution}{exo:b1:acetals-protection:8}
डाइऑल के साथ एक अणु दो का स्थान लेता है: अभिक्रिया अणुओं की संख्या
नहीं घटाती (एक ऐल्डिहाइड और एक डाइऑल एक \omterm{def:b1:acetals-protection:hemiacetal}{ऐसीटैल} और एक
जल देते हैं), और दूसरा OH अभिक्रिया करने वाले कार्बन के पास बना रहता है, जिससे
वलय आसानी से बंद होता है। दोनों बातें चक्रीय \omterm{def:b1:acetals-protection:hemiacetal}{ऐसीटैल} को अधिक अनुकूल बनाती हैं।
\end{solution}

\begin{solution}{exo:b1:acetals-protection:9}
अम्ल \omterm{def:b1:organomagnesium:organometallic}{ग्रीन्यार अभिकर्मक} को नष्ट कर देगा (प्रोटॉन स्थानांतरण); जल और डाइऑल के
अंश भी हटाने होंगे।
\end{solution}

\begin{solution}{exo:b1:acetals-protection:10}
1. 4-क्लोरोब्यूटेन-2-ओन के कीटोन का चक्रीय \omterm{def:b1:acetals-protection:hemiacetal}{ऐसीटैल} के रूप में \omterm{def:b1:acetals-protection:protecting-group}{रक्षण} कीजिए
(एथेन-1,2-डाइऑल, अम्ल \omterm{def:b1:elementary-steps:catalyst}{उत्प्रेरक}, डीन--स्टार्क)। 2. शुष्क ईथर में मैग्नीशियम:
रक्षित क्लोराइड का \omterm{def:b1:organomagnesium:organometallic}{ग्रीन्यार अभिकर्मक}, जो अब अपने ही
कार्बोनिल से नष्ट नहीं होता। 3. प्रोपेनोन मिलाइए; जलीय अमोनियम क्लोराइड से उपचार:
तृतीयक ऐल्कोहॉल। 4. तनु जलीय अम्ल: \omterm{def:b1:acetals-protection:protecting-group}{विरक्षण} से
5-हाइड्रॉक्सी-5-मेथिलहेक्सेन-2-ओन।
\end{solution}

\begin{solution}{exo:b1:acetals-protection:11}
$Q = [\text{ऐसीटैल}][\ce{H2O}]/([\text{ऐल्डिहाइड}][\ce{CH3OH}]^2)$: मेथेनॉल का बड़ा
आधिक्य, जो हर में वर्गित है, $Q < K$ बनाए रखता है जब तक थोड़ा ही
ऐल्डिहाइड न बचे। जल हटाना अंश पर कार्य करता है; मेथेनॉल को
विलायक के रूप में लेना सरल है पर उसके लिए आसानी से पृथक होने वाला उत्पाद चाहिए।
\end{solution}

\begin{solution}{exo:b1:acetals-protection:12}
केवल \omterm{def:b1:acetals-protection:hemiacetal}{हेमीऐसीटैल} का OH जल के रूप में निकलकर ऐसा धनायन दे सकता है जो
पड़ोसी वलय-ऑक्सीजन द्वारा स्थायी हो (\ce{C=O+}); अन्य OH समूहों को
साधारण कार्बनों से निकलना पड़ता, जिससे अस्थायी धनायन बनते। मेथेनॉल
स्थायी धनायन से जुड़ता है: एक मेथिल \omterm{def:b1:acetals-protection:hemiacetal}{ऐसीटैल}।
\end{solution}

\begin{solution}{pb:b1:acetals-protection:1}
\textbf{1.} \ce{BrMgCH2CH2CH2CHO}।
\textbf{2.} उसका \omterm{def:b1:electronic-effects:intermediates}{कार्बऋणायन} जैसा कार्बन किसी अन्य अणु के (या स्वयं के) ऐल्डिहाइड पर
आक्रमण करेगा: \ce{RMgBr}, $\mathrm{R'CHO}$ से जुड़कर एक \omterm{def:b1:alcohol-activation:alkoxide}{ऐल्कॉक्साइड} देता है।
\textbf{3.} ऐल्डिहाइड, \omterm{def:b1:organomagnesium:organometallic}{ग्रीन्यार अभिकर्मक} के विरुद्ध।
\textbf{4.} \omterm{def:b1:acetals-protection:hemiacetal}{ऐसीटैल} \omterm{def:b1:organomagnesium:organometallic}{ग्रीन्यार अभिकर्मकों} और क्षारकों के प्रति स्थायी हैं और
जलीय अम्ल से हटाए जाते हैं।
\textbf{5.} वह अधिक पूर्णता से बनता है (ऐल्कोहॉल के दो अणुओं के लिए एक डाइऑल)
और उसका जल-अपघटन आसान है।
\textbf{6.} $15.09/150.9 = \qty{0.1000}{mol}$।
\textbf{7.} $1.20 \times 0.1000 \times 62.0 = \qty{7.44}{g}$।
\textbf{8.} \ce{BrCH2CH2CH2CHO + HOCH2CH2OH <=> C6H11BrO2 + H2O}
(2-(3-ब्रोमोप्रोपिल)-1,3-डाइऑक्सोलेन)।
\textbf{9.} वह कार्बोनिल को सक्रिय करता है और दूसरे चरण में जल की हानि
संभव बनाता है; वह पुनर्जनित होता है।
\textbf{10.} टॉलूईन और जल साथ-साथ आसवित होते हैं, संघनित होकर
संग्राहक में गिरते हैं; जल नीचे बैठकर रुका रहता है, टॉलूईन उफनकर लौट जाता है।
\textbf{11.} $Q$ के अंश में $[\ce{H2O}]$ है; जल हटाने से
$Q < K$ बना रहता है, और अभिक्रिया आगे बढ़ती है।
\textbf{12.} जल अधिक घना है और टॉलूईन के साथ नहीं मिलता: वह तली में
एकत्रित होता है; केवल ऊपरी परत वापसी भुजा तक पहुँचती है।
\textbf{13.} जब अपेक्षित आयतन पर पहुँचकर जल का एकत्रित होना रुक जाए।
\textbf{14.} $0.1000 \times 194.9 = \qty{19.5}{g}$।
\textbf{15.} \ce{C6H11BrO2 + Mg -> C6H11BrMgO2} (शुष्क ईथर)।
\textbf{16.} अभिकर्मक \ce{(CH3)2CO} से जुड़ता है: रक्षित शृंखला पर
\omterm{def:b1:alcohol-activation:alkoxide}{ऐल्कॉक्साइड} \ce{(CH3)2C(O^-)CH2CH2CH2-}, \ce{MgBr+} के साथ।
\textbf{17.} अम्ल \omterm{def:b1:acetals-protection:protecting-group}{रक्षी समूह} को बहुत जल्दी हटा देगा और
तृतीयक ऐल्कोहॉल का \omterm{def:b1:alcohol-activation:dehydration}{निर्जलीकरण} कर सकता है।
\textbf{18.} $0.0700 \times 174.0 = \qty{12.2}{g}$।
\textbf{19.} \omterm{def:b1:acetals-protection:hemiacetal}{ऐसीटैल} + \ce{H2O} $\to$ ऐल्डिहाइड + एथेन-1,2-डाइऑल
(अम्ल \omterm{def:b1:elementary-steps:catalyst}{उत्प्रेरक})।
\textbf{20.} उसके जल खोने के लिए \omterm{def:b1:predominant-reaction:strong-acid}{प्रबल अम्ल} और ऊष्मा चाहिए; मृदु, ठंडा
जलीय अम्ल \omterm{def:b1:acetals-protection:hemiacetal}{ऐसीटैल} का जल-अपघटन कहीं अधिक तेज़ी से करता है।
\textbf{21.} कार्बन 5 का OH ऐल्डिहाइड कार्बन 1 से जुड़ता है: एक छह-सदस्यीय
वलय (पाँच कार्बन और एक ऑक्सीजन), जिसमें वलय-ऑक्सीजन के बगल वाले एक कार्बन पर दो मेथिल
समूह हैं और उसके बगल वाले दूसरे कार्बन पर एक OH।
\textbf{22.} $0.0700 \times 0.90 = \qty{0.0630}{mol}$, \qty{8.19}{g}।
\textbf{23.} $0.0630/0.1000 = \qty{63}{\percent}$ (\omterm{def:b1:acetals-protection:protecting-group}{रक्षण} को
पूर्ण मानकर)।
\textbf{24.} $0.1000 \times 18.0 = \qty{1.80}{g}$।
\textbf{25.} $1.80/0.997 =$ जाल में \textbf{\qty{1.8}{mL} जल}।
\end{solution}
